\documentclass{article}
\usepackage{amsmath, amssymb, amsthm}

\title{ACM India Summer School on Symmetric Key Cryptography \\ Day 2: Negligible Functions, PRGs, and PRFs Solutions}
\author{Sourav Sharma}
\date{June 9, 2026}

\newtheorem{claim}{Claim}
\newtheorem{theorem}{Theorem}

\begin{document}

\maketitle

\section{Negligible Functions}

A function $g(k)$ is negligible (denoted $g(k) \in \text{negl}(k)$) if for every constant $c > 0$, there exists a $k' \in \mathbb{N}$ such that for all $k \geq k'$:
\[
|g(k)| < \frac{1}{k^c}
\]

\begin{claim}
Let $\text{negl}_1$ and $\text{negl}_2$ be two negligible functions. Then $g(k) = \text{negl}_1(k) + \text{negl}_2(k)$ is also negligible.
\end{claim}
\begin{proof}
We need to prove that for any arbitrary constant $c > 0$, there exists a $k'$ such that for all $k \geq k'$, $\text{negl}_1(k) + \text{negl}_2(k) < \frac{1}{k^c}$.

Let $c' = c + 1$. Since $\text{negl}_1$ and $\text{negl}_2$ are negligible:
\begin{itemize}
    \item There exists $k_1$ such that for all $k \geq k_1$, $\text{negl}_1(k) < \frac{1}{2k^c}$ (by considering the polynomial $2k^c$).
    \item There exists $k_2$ such that for all $k \geq k_2$, $\text{negl}_2(k) < \frac{1}{2k^c}$.
\end{itemize}
Let $k' = \max\{k_1, k_2\}$. For all $k \geq k'$:
\[
\text{negl}_1(k) + \text{negl}_2(k) < \frac{1}{2k^c} + \frac{1}{2k^c} = \frac{1}{k^c}
\]
Thus, the sum is negligible.
\end{proof}

\begin{claim}
A negligible function stays negligible even after polynomial multiplication: $\text{poly}(k) \cdot \text{negl}(k) = \text{negl}(k)$.
\end{claim}
\begin{proof}
Let $p(k) = k^a$ for some constant $a > 0$ represent the polynomial, and $g(k)$ be a negligible function. We want to show that $p(k) \cdot g(k)$ is negligible.
For any constant $c > 0$, let $c' = c + a$. Since $g(k)$ is negligible, for the constant $c'$, there exists $k'$ such that for all $k \geq k'$:
\[
g(k) < \frac{1}{k^{c'}} = \frac{1}{k^{c+a}}
\]
Multiplying both sides by $p(k) = k^a$:
\[
p(k) \cdot g(k) < k^a \cdot \frac{1}{k^{c+a}} = \frac{1}{k^c}
\]
Thus, $p(k) \cdot g(k)$ is negligible.
\end{proof}

\begin{claim}
If $f(k)$ is not a negligible function and $g(k)$ is a negligible function, then $h(k) = f(k) - g(k)$ is not a negligible function.
\end{claim}
\begin{proof}
We prove this by contradiction. Assume that $h(k) = f(k) - g(k)$ is a negligible function. 
Since the sum of two negligible functions is negligible (by Claim 1), and both $h(k)$ and $g(k)$ are negligible under our assumption, it follows that:
\[
f(k) = h(k) + g(k)
\]
must also be a negligible function.
This directly contradicts the given assumption that $f(k)$ is not a negligible function. Therefore, the assumption that $f(k) - g(k)$ is negligible must be false. Hence, $f(k) - g(k)$ is not negligible.
\end{proof}

\section{Pseudorandom Generators (PRGs)}

\begin{claim}
Let $G_1, G_2 : \{0, 1\}^n \to \{0, 1\}^{2n}$ be two secure PRGs with $G_1 \neq G_2$. Define $G(s) = G_1(s) \oplus G_2(s)$. Is $G$ necessarily a PRG?
\end{claim}
\begin{proof}
No, $G$ is not necessarily a PRG. We show a counterexample.
Let $G_1 : \{0,1\}^n \to \{0,1\}^{2n}$ be a secure PRG. Define $G_2(s) = G_1(s) \oplus 1^{2n}$. Since $G_1$ is a secure PRG, its XOR with a constant string $1^{2n}$ is also a secure PRG. However, computing $G(s)$ gives:
\[
G(s) = G_1(s) \oplus G_2(s) = G_1(s) \oplus G_1(s) \oplus 1^{2n} = 1^{2n}
\]
The output of $G(s)$ is always the constant string $1^{2n}$. A distinguisher $D(w)$ that outputs $1$ if $w = 1^{2n}$ and $0$ otherwise distinguishes $G(s)$ from the uniform distribution $U_{2n}$ with probability $1 - 2^{-2n}$, which is non-negligible. Thus, $G$ is not a secure PRG.
\end{proof}

\begin{claim}
Let $G$ be a secure PRG with expansion factor $\ell(n) > 2n$. Say whether $G_1$ is necessarily a PRG:
\begin{enumerate}
    \item $G_1(s) = G(s_1 \dots s_{\lfloor n/2 \rfloor})$
    \item $G_1(s) = G(0^{|s|} \parallel s)$
    \item $G_1(s) = G(s) \parallel G(s+1)$
\end{enumerate}
\end{claim}
\begin{proof}[Solutions]
\hfill
\begin{enumerate}
    \item \textbf{No.} $G_1$ does not satisfy the expansion requirement. If $\ell(n) = 2n + 1$, then for the input of length $n$, the output of $G_1(s)$ has length $\ell(\lfloor n/2 \rfloor) = 2\lfloor n/2 \rfloor + 1 \leq n + 1$, which is not strictly greater than the input size $n$ for all cases (e.g. for odd $n$).
    
    \item \textbf{No.} We can construct a counterexample. Let $\hat{G} : \{0,1\}^{2n} \to \{0,1\}^{\ell(2n)}$ be a secure PRG. Define $G(s_1)$ as:
    \[
    G(s_1) = \begin{cases}
    0^{\ell(2n)} & \text{if the first } n \text{ bits of } s_1 \text{ are all zeros} \\
    \hat{G}(s_1) & \text{otherwise}
    \end{cases}
    \]
    Since the probability of a uniform $s_1$ starting with $0^n$ is $2^{-n}$ (which is negligible), $G$ is still a secure PRG. But if we define $G_1(s) = G(0^{|s|} \parallel s)$, the first half of its seed is always $0^{|s|}$, meaning $G_1(s)$ always outputs $0^{\ell(2n)}$. This is trivially distinguishable from a random distribution.
    
    \item \textbf{No.} We show a counterexample. Let $G' : \{0,1\}^n \to \{0,1\}^{\ell(n)-1}$ be a secure PRG. We construct a secure PRG $G : \{0,1\}^n \to \{0,1\}^{\ell(n)}$ as follows:
    \[
    G(s) = \begin{cases}
    G'(s) \parallel 0 & \text{if } s \text{ is even (i.e. the LSB of } s \text{ is } 0\text{)} \\
    G'(s-1) \parallel 1 & \text{if } s \text{ is odd (i.e. the LSB of } s \text{ is } 1\text{)}
    \end{cases}
    \]
    Since LSB of a uniform seed $s$ is $0$ or $1$ with probability $1/2$, the last bit of $G(s)$ is uniform, and $G(s)$ is a secure PRG.
    
    Now, let's analyze $G_1(s) = G(s) \parallel G(s+1)$ where $s+1$ is addition modulo $2^n$:
    \begin{itemize}
        \item If $s$ is even, $s+1$ is odd. Then:
        \[
        G_1(s) = [G'(s) \parallel 0] \parallel [G'((s+1)-1) \parallel 1] = G'(s) \parallel 0 \parallel G'(s) \parallel 1
        \]
        Here, the last bit of the first half is $0$, and the last bit of the second half is $1$.
        \item If $s$ is odd, $s+1$ is even. Then:
        \[
        G_1(s) = [G'(s-1) \parallel 1] \parallel [G'(s+1) \parallel 0]
        \]
        Here, the last bit of the first half is $1$, and the last bit of the second half is $0$.
    \end{itemize}
    In both cases, the last bit of the first half of $G_1(s)$ and the last bit of the second half are always different (i.e. their XOR is always $1$). A distinguisher can check if the XOR of these two bits is $1$. For $G_1(s)$ it is always $1$ (probability $1$), whereas for a uniform random string it is $1$ with probability $1/2$. This yields a distinguishing advantage of $1/2$, proving $G_1$ is not a secure PRG.
\end{enumerate}
\end{proof}

\section{Pseudorandom Functions (PRFs)}

\begin{claim}
Consider $F_{A,b}(x) = Ax + b \pmod 2$, where $A$ is an $n \times n$ boolean matrix and $b$ is an $n$-bit vector. Show that $F$ is not a secure PRF.
\end{claim}
\begin{proof}
A secure PRF is computationally indistinguishable from a truly random function. However, the construction $F_{A,b}(x) = Ax + b \pmod 2$ is linear (affine). It satisfies the following relation for any inputs $x, y \in \{0, 1\}^n$:
\[
F_{A,b}(x) \oplus F_{A,b}(y) \oplus F_{A,b}(x \oplus y) = (Ax + b) \oplus (Ay + b) \oplus (A(x \oplus y) + b)
\]
\[
= Ax \oplus b \oplus Ay \oplus b \oplus Ax \oplus Ay \oplus b = b
\]
Thus, the XOR sum $F_{A,b}(x) \oplus F_{A,b}(y) \oplus F_{A,b}(x \oplus y)$ is always constant and equals the vector $b$ for all $x, y$.

We construct a distinguisher $D$ that queries the oracle $O$ on:
\begin{enumerate}
    \item $u, v$ and computes $t_1 = O(u) \oplus O(v) \oplus O(u \oplus v)$.
    \item $x, y$ (distinct from $u, v$) and computes $t_2 = O(x) \oplus O(y) \oplus O(x \oplus y)$.
\end{enumerate}
If $t_1 = t_2$, the distinguisher outputs $1$ (guessing $O$ is $F_{A,b}$); otherwise it outputs $0$.
\begin{itemize}
    \item If $O = F_{A,b}$, then $t_1 = t_2 = b$ always, so $D$ outputs $1$ with probability $1$.
    \item If $O$ is a truly random function, $t_1$ and $t_2$ are independent uniform random $n$-bit strings, so $t_1 = t_2$ with probability $2^{-n}$.
\end{itemize}
The distinguishing advantage is $1 - 2^{-n}$, which is non-negligible. Thus, $F_{A,b}$ is not a secure PRF.
\end{proof}

\begin{claim}
Let $F$ be a length-preserving pseudorandom function. Show that the keyed function $F': \{0, 1\}^{n-1} \to \{0, 1\}^{2n}$ defined by $F'_k(x) = F_k(0 \parallel x) \parallel F_k(1 \parallel x)$ is a pseudorandom function.
\end{claim}
\begin{proof}
We prove security of $F'$ using a hybrid argument. We assume for contradiction that there exists a PPT distinguisher $D$ that distinguishes $F'_k(x)$ from a truly random function $R: \{0, 1\}^{n-1} \to \{0, 1\}^{2n}$ with non-negligible advantage.

Define the following hybrid distributions for a fixed input $x$:
\begin{itemize}
    \item $H_0 = F_k(0 \parallel x) \parallel F_k(1 \parallel x)$ (which is the output of $F'_k(x)$)
    \item $H_1 = U_n \parallel F_k(1 \parallel x)$ (where the first half is replaced by a uniform $n$-bit random string $U_n$)
    \item $H_2 = U_n \parallel U_n$ (which is a uniformly random $2n$-bit string)
\end{itemize}

If $D$ distinguishes $H_0$ from $H_2$ with non-negligible advantage, then by the hybrid argument, $D$ must distinguish $H_0$ from $H_1$, or $H_1$ from $H_2$ (or both) with non-negligible advantage.

\paragraph{Case 1: Distinguishing $H_0$ from $H_1$.}
Suppose $| \Pr[D(H_0) = 1] - \Pr[D(H_1) = 1] |$ is non-negligible. We construct a PPT distinguisher $D_1$ against the PRF $F$. Given oracle access to a function $\mathcal{O}$ (which is either $F_k$ or a truly random function), $D_1$ does the following:
\begin{enumerate}
    \item Query $\mathcal{O}$ on $0 \parallel x$ to obtain $y$.
    \item Independently compute $F_k(1 \parallel x)$.
    \item Query $D$ on the input $y \parallel F_k(1 \parallel x)$.
    \item Output whatever $D$ outputs.
\end{enumerate}
If $\mathcal{O} = F_k$, the input to $D$ is distributed exactly as $H_0$. If $\mathcal{O}$ is a truly random function, the input to $D$ is distributed as $H_1$. Thus, $D_1$ distinguishes $F$ from a random function with non-negligible advantage, contradicting the security of $F$ as a PRF.

\paragraph{Case 2: Distinguishing $H_1$ from $H_2$.}
Suppose $| \Pr[D(H_1) = 1] - \Pr[D(H_2) = 1] |$ is non-negligible. We construct a PPT distinguisher $D_2$ against the PRF $F$ similarly. Given oracle access to $\mathcal{O}$:
\begin{enumerate}
    \item Query $\mathcal{O}$ on $1 \parallel x$ to obtain $y$.
    \item Choose a uniform random string $r \leftarrow \{0, 1\}^n$.
    \item Query $D$ on the input $r \parallel y$.
    \item Output whatever $D$ outputs.
\end{enumerate}
If $\mathcal{O} = F_k$, the input to $D$ is distributed exactly as $H_1$. If $\mathcal{O}$ is a truly random function, the input to $D$ is distributed as $H_2$. Thus, $D_2$ distinguishes $F$ from a random function with non-negligible advantage, contradicting the security of $F$.

Since both cases lead to contradictions, the transition from $H_0$ to $H_2$ cannot be distinguished with non-negligible advantage. Therefore, $H_0 \approx_c H_2$, and $F'_k(x)$ is a pseudorandom function.
\end{proof}

\end{document}
